\section{\arcto{}: A HIP Compression Runtime for AMD GPUs}
\label{sec:arcto-tool}

\arcto{} is a HIP-based runtime for GPU data compression on AMD
platforms. It exposes a batched chunked execution model derived from
nvCOMP-style APIs and currently supports three lossless byte-level
codecs (LZ4, Snappy, Cascaded) together with an experimental
error-bounded backend based on ZFP for floating-point scientific data.
The runtime is compiled with \texttt{hipcc} for explicit AMD GFX
targets covering both CDNA (\texttt{gfx906}, \texttt{gfx90a},
\texttt{gfx942}) and RDNA3 (\texttt{gfx1100}) generations, so the
same source tree services the four AMD parts used in the evaluation.

The starting point for the runtime is a HIP port of nvCOMP-style
batched compression kernels~\cite{anon_iccsa2026}. The port follows
the usual CUDA-to-HIP path of mechanical API translation, replacement
of CUDA-specific primitives with HIP or hipCUB equivalents, CMake
adaptation for \texttt{hipcc}, and correctness validation through
lossless round-trip tests. That baseline established that the kernels
themselves can be moved to AMD with bit-exact reconstruction, but it
also showed that on host-resident inputs, PCIe transfers dominate
end-to-end time even when the compression kernels achieve high
throughput. \arcto{} is the runtime layer that sits on top of these
ported kernels and addresses the two pieces of the pipeline that the
characterization left underexploited: the kernel-visible chunk
granularity, and the host-device transfer policy. A third extension,
the ZFP backend, broadens the runtime to data classes for which
lossless ratios are too small to justify compression overhead.

\subsection{Compression backends and execution model}
\label{sec:arcto-execution-model}

The three lossless codecs follow the same batched abstraction. Each
input buffer is split into independent chunks, and each chunk is
compressed or decompressed independently on the GPU. This exposes
many work items, avoids inter-chunk dependencies, and makes the
kernels naturally suited to massively parallel execution. LZ4 uses
a private hash table per chunk for match finding, which makes it
sensitive to memory hierarchy behavior. Snappy maps cooperating
wavefronts to chunks rather than threads, so its scaling differs
from LZ4. Cascaded applies a structured pipeline of delta encoding,
run-length encoding, and bit packing, which is more regular than the
dictionary-style codecs and tends to be more stable across input
distributions. The ZFP backend targets floating-point data and
operates on small blocks rather than byte chunks; it is described in
Section~\ref{sec:arcto-zfp}.

A central design choice of \arcto{} is that the granularity at which
the kernels see the input is not the same as the granularity at which
the host stages data to the device. The runtime therefore treats two
quantities separately: a kernel-visible chunk size $C_{\text{ker}}$,
which controls the amount of independent work exposed to the
compression kernels, and a pinned aggregation window $W_{\text{agg}}$,
which controls how much host data is staged through pinned memory
before transfer to the GPU. Sections~\ref{sec:optim-chunk-size}
and~\ref{sec:optim-adaptive-tiled} develop the two parameters in
detail; what matters here is that \arcto{} can keep $C_{\text{ker}}$
small for kernel parallelism while grouping many such chunks into a
larger transfer-efficient window. The kernels themselves are
unchanged: only the runtime policy around them is.

\subsection{Host execution paths and instrumentation}
\label{sec:arcto-host-paths}

The runtime offers three host execution paths. The pageable path is
the baseline: the input resides in ordinary host memory and HIP
handles the H2D and D2H transfers directly. It is simple and
compatible with existing applications, but it leaves substantial
end-to-end time on the table for inputs of a few hundred MiB and
above. The full pinned path allocates a pinned host buffer the same
size as the input, copies the pageable input into it, and transfers
from pinned memory to the GPU. This reduces H2D and D2H times sharply
but introduces an allocation cost, a host-to-host staging cost, and a
peak pinned-memory footprint that scales with the input. The adaptive
pinned aggregation path resolves this tradeoff: it allocates a
fixed-size pinned aggregation window once and reuses it across
successive portions of the input, so the peak pinned footprint is
bounded by $W_{\text{agg}}$ rather than by $L$. The adaptive path
changes only the transfer policy, not the compression kernels.

Beyond the three host paths, \arcto{} is designed to also support a
GPU-resident workflow in which the data is produced directly by an
upstream GPU kernel (for example, a finite-difference wave-propagation
solver) and never lives in host memory before compression. In that
case the H2D phase and the host-to-host staging disappear, and only
the compression kernel time and the compressed D2H volume remain on
the critical path. The present paper does not develop GPU-resident
execution as a separate contribution, but
Section~\ref{sec:evaluation} reports a preliminary in-situ
integration of \arcto{} into a HIP-ported seismic propagator as a
validation case.

Because the bottleneck moves between phases depending on which path
is active, the runtime exposes a detailed timing breakdown:
allocation time $T_{\text{alloc}}$, host-to-host staging time
$T_{\text{h2h}}$, host-to-device transfer time $T_{\text{h2d}}$,
kernel time $T_{\text{kernel}}$, device-to-host transfer time
$T_{\text{d2h}}$, and total time $T_{\text{total}}$, together with
peak pinned bytes, number of adaptive windows, and compression ratio.
This breakdown is necessary because pageable, pinned, and adaptive
modes look very different depending on whether one measures only
kernel execution, only transfers, or the full host-side pipeline; the
evaluation in Section~\ref{sec:evaluation} relies on this
decomposition to avoid the common pitfall of reporting pinned-memory
speedups while omitting the allocation and staging costs that the
pinned path actually pays.

\subsection{Error-bounded compression with ZFP}
\label{sec:arcto-zfp}

The lossless codecs described above achieve only marginal compression
ratios on the floating-point scientific data that motivates this work.
On the TTI seismic wavefields used in the evaluation, LZ4, Snappy, and
Cascaded reach ratios of roughly
1.03--1.04$\times$~\cite{anon_iccsa2026}, which is rarely enough to
justify compression overhead when the objective is reducing
checkpoint or migration-storage volume. \arcto{} therefore includes
an experimental ZFP backend for floating-point data, where bounded
numerical error can be traded for much larger reductions.

The backend exposes the three canonical ZFP modes: a fixed-rate mode,
in which the user specifies a target number of compressed bits per
value; a fixed-accuracy mode, in which the user specifies an absolute
error tolerance; and a fixed-precision mode, in which the user
specifies the number of retained bit planes. For seismic wave
propagation, the fixed-rate mode is particularly attractive because
it gives a predictable per-checkpoint storage and transfer volume:
a single-precision wavefield uses 32 bits per sample uncompressed,
and a fixed-rate ZFP configuration can reduce that to a chosen
number of bits per sample while bounding the reconstruction error.
The kernels follow the same batched style as the lossless codecs:
input data is partitioned into independent blocks, compressed on the
GPU, and transferred to the host in compressed form. Prior work on
reverse-time migration has shown that selected lossy wavefield
configurations can preserve the quality of the final migrated image
while reducing wavefield storage
requirements~\cite{barbosa2023rtmcompression}, which is the
motivating use case for the ZFP path in \arcto{}.

\subsection{Build and validation}
\label{sec:arcto-build}

\arcto{} is built with CMake and \texttt{hipcc}, with the target
architectures selected through \texttt{CMAKE\_HIP\_ARCHITECTURES}.
The wavefront width is set at compile time, so CDNA builds use
64-thread wavefronts and the RDNA3 build uses wave32; this avoids a
rewrite of the warp-level logic inside the compression kernels.
Correctness for the lossless codecs is verified through a round-trip
test that requires bit-exact equality
$\operatorname{decompress}(\operatorname{compress}(x)) = x$ for every
combination of algorithm, dataset, architecture, and chunk size used
in the evaluation. ZFP modes are validated mode-dependently against
the configured error bound, with absolute and relative error
statistics reported alongside compression ratio and throughput. For
reproducibility, the benchmark environment uses a containerized
ROCm software stack and writes platform metadata, selected chunk
size, transfer mode, the full timing breakdown, and the runtime
configuration into the per-run CSV output.
